\begin{abstract}
% 1. State the problem
General-purpose database engines apply the same fixed rewrite rules to every workload.
So they miss the workload-specific rewrites and phase ordering that can yield order-of-magnitude speedups.
% 2. Say why it's an interesting problem
LLM-based rewriters discover such rewrites via whole-query rewriting, but they pay high costs in terms of LLM tokens and inference time for every query or template. These rewrites do not readily transfer to new queries. Furthermore, engine synthesis generates a bespoke engine per workload, which requires re-synthesis when queries change (costing hundreds of dollars for a few queries), and it does not reuse the execution engines of mature systems.
We introduce {\sys}, a system that teaches existing database engines new rewrites.
{\sys} runs a pluggable LLM rewriter 
offline on a sample of each workload for discovering rewrites with moderate LLM costs.
It generalizes each rewrite into a rule that transfers across workloads.
Such rules are only accepted if they produce equivalent results (using SMT solvers or result-equivalence tests).
At query time, {\sys} matches rules on shared plan fragments of logical plans, rewrites a query without any LLM calls, and reuses the host execution engine.
% 4. Say what follows from your solution
We evaluate {\sys} on nine workloads over PostgreSQL, MySQL, and DuckDB, with all rewrites verified correct.
On Stats-CEB and PostgreSQL, the distilled rule sets yield workload speedups of $38.7\times$, compared to $12.9\times$ for ReSequel, a state-of-the-art LLM rewriter.
Rules distilled from JOB and transferred to Stats-CEB yield a speedup of $28.3\times$.
On CEB and DuckDB with multi-threading, {\sys} outperforms BespokeOLAP, an LLM-synthesized engine, by $2.05$--$2.25\times$ without any LLM costs.
\end{abstract}
